\documentclass[9pt]{extarticle}
\usepackage[a4paper,landscape,margin=10mm]{geometry}
\usepackage[french]{babel}
\usepackage[most]{tcolorbox}
\usepackage[os=win]{menukeys}
\usepackage{enumitem}
\usepackage{environ}
\usepackage{fontawesome}
\usepackage{fontspec}
\usepackage{graphicx}
\usepackage{listings}
\usepackage{microtype}
\usepackage{multicol}
\usepackage{ragged2e}
\usepackage{tabularx}
\usepackage{titlesec}
\usepackage{xparse}

\setlength{\columnsep}{5mm}
\setlength{\columnseprule}{0.4pt}

\setlist{nosep,leftmargin=1.1em}
\setlength{\parindent}{0pt}
\setlength{\parskip}{2pt}

% \setmainfont{TeX Gyre Heros}
% \setsansfont{TeX Gyre Heros}
% \setmonofont{Inconsolata}
\renewcommand{\familydefault}{\sfdefault}

\lstset{
  basicstyle=\ttfamily\footnotesize,
  columns=fullflexible,
  keepspaces=true,
  upquote=true,
  breaklines=true,
  literate=*%
    {[f]}{{\raisebox{-0.15ex}{\scriptsize\faFileO}}}3
    {[fc]}{{\raisebox{-0.15ex}{\scriptsize\faFileCodeO}}}4
    {[fa]}{{\raisebox{-0.15ex}{\scriptsize\faFileArchiveO}}}4
    {[d]}{{\raisebox{-0.15ex}{\scriptsize\faFolderO}}}3
    {[p]}{{\raisebox{-0.15ex}{\scriptsize\faFilePowerpointO}}}2
    {[t]}{{\raisebox{-0.15ex}{\scriptsize\faHashtag}}}2
    {uu}{{test}}2
}

\NewDocumentCommand{\code}{v}{%
  \tcbox[
    on line,
    colframe=black!60, colback=black!5,
    bottom=1.5pt, left=1.5pt, right=1.5pt, top=1.5pt,
    boxrule=0.4pt, arc=2pt, boxsep=0.2pt,
    fontupper=\ttfamily\footnotesize,
    before upper=\rule{0pt}{0.7em},
    tcbox raise base
  ]{\lstinline!#1!}%
}

\NewDocumentCommand{\codx}{m}{%
  \tcbox[
    on line,
    colframe=black!60, colback=black!5,
    bottom=1.5pt, left=1.5pt, right=1.5pt, top=1.5pt,
    boxrule=0.4pt, arc=2pt, boxsep=0.2pt,
    fontupper=\ttfamily\footnotesize,
    before upper=\rule{0pt}{0.7em},
    tcbox raise base
  ]{#1}%
}

\renewmenumacro{\keys}[+]{roundedkeys}

\let\oldkeys\keys
\renewcommand{\keys}[1]{\small\oldkeys{#1}\normalsize}

% Tailles des titres (s'applique aussi aux versions étoilées *)
\titleformat{\section}
  {\sffamily\bfseries\fontsize{12}{15}\selectfont} % ajoute  si tu veux en gras
  {}{0pt}{}

\titleformat{\subsection}
  {\sffamily\bfseries\fontsize{10}{12}\selectfont}
  {}{0pt}{}

\titleformat{\subsubsection}
  {\sffamily\bfseries\fontsize{8}{10}\selectfont}
  {}{0pt}{}

% Espacements (les tiens, inchangés)
\titlespacing*{\section}{0pt}{0.1ex plus .1ex minus 0.1ex}{0.5ex plus .1ex minus 0.1ex}
\titlespacing*{\subsection}{0pt}{0.1ex plus .1ex minus 0.1ex}{0.5ex plus .1ex minus 0.1ex}
\titlespacing*{\subsubsection}{0pt}{0.1ex plus .1ex minus 0.1ex}{0.5ex plus .1ex minus 0.1ex}


\pagenumbering{gobble} % pas de numéros de page


\newlength\mybaselinestretch
\mybaselinestretch=0pt plus 0.02pt\relax
\addtolength{\baselineskip}{\mybaselinestretch}

\setlength\parindent{0pt}
\setlength\tabcolsep{1.5pt}
\setlength{\columnseprule}{0.4pt}

\input{revision}

\newenvironment{cmdblock}{%
  \par\setlength{\parindent}{0pt}\setlength{\parskip}{0pt}%
  \RaggedRight%
}{\par}

\newcommand{\cmditem}[2]{%
  \noindent\inlinecode{#1}\quad #2\par
}

% indentation supplémentaire pour les sous-lignes (options)
\newlength{\cmdoptindent}
\setlength{\cmdoptindent}{1em} % ajuste au goût

\newcommand{\cmdopt}[1]{%
  \hspace*{\cmdoptindent}%
  \begin{minipage}[t]{\dimexpr\linewidth-\cmdoptindent\relax}
    \RaggedRight \itshape #1%
  \end{minipage}\par
}

\newcommand{\codefa}[2][]{%
  \raisebox{-0.15ex}{\faIcon[#1]{#2}}%
}

\clubpenalty=0            % orphelines (début de page)
\widowpenalty=0           % veuves (fin de page)
\displaywidowpenalty=0

\setlength{\premulticols}{0pt} % pas d'espace avant/après l'environnement
\setlength{\postmulticols}{0pt}
\setlength{\multicolsep}{0pt}

\begin{document}

\raggedcolumns
\begin{multicols}{4}
    \raggedbottom

    \begin{tabularx}{\columnwidth}{lX}
        \raisebox{-2em}{\includegraphics[height=5em]{assets/linux-logo.pdf}} &
        \begin{raggedright}
            {\Large \bf Linux / Unix} \\
            HEIG-VD -- version \revision \ -- \revisiondate \\
        \end{raggedright}
    \end{tabularx}
    \begin{raggedright}
        Carte de référence des commandes et utilitaires Unix-like courants. Les marqueurs suivants sont utilisés : \code{[f]} : nom de fichier,
        \code{[d]} : nom de répertoire,
        \code{[p]} : commande,
        \code{"s"} : chaîne de caractères
    \end{raggedright}

    \section*{Raccourcis clavier Bash}
    \subsection*{Mouvements}

    \begin{tabularx}{\columnwidth}{lX}
        \keys{\ctrl + B} & recule d'un caractère (\keys{\leftarrow})   \\
        \keys{\ctrl + F} & avance d'un caractère (\keys{\rightarrow})  \\
        \keys{\ctrl + A} & début de ligne (\keys{Home})                \\
        \keys{\ctrl + E} & fin de ligne (\keys{End})                   \\
        \keys{\Alt + B}  & recule d'un mot (\keys{\Alt + \leftarrow})  \\
        \keys{\Alt + F}  & avance d'un mot (\keys{\Alt + \rightarrow}) \\
    \end{tabularx}

    \subsection*{Édition}

    \begin{tabularx}{\columnwidth}{lX}
        \keys{\ctrl + H} & comme backspace (\keys{\backspace})    \\
        \keys{\ctrl + D} & comme delete (\keys{\del})             \\
        \keys{\ctrl + U} & coupe du début au curseur              \\
        \keys{\ctrl + K} & coupe du curseur à la fin              \\
        \keys{\ctrl + W} & coupe le mot avant le curseur          \\
        \keys{\ctrl + Y} & colle le dernier texte coupé           \\
        \keys{\ctrl + T} & swap le char courant avec le précédent \\
        \keys{\Alt + D}  & supprime le mot après                  \\
        \keys{\Alt + .}  & insère le dernier argument             \\
    \end{tabularx}

    \subsection*{Exécution / complétion}

    \begin{tabularx}{\columnwidth}{lX}
        \multicolumn{2}{l}{\keys{\ctrl + J}, \keys{\ctrl + M} exécuter (\keys{\return})} \\
        \keys{\ctrl + I} & complétion (équivalent \keys{\tab})                           \\
        \keys{\tab}{}    & auto-complétion                                               \\
    \end{tabularx}

    \subsection*{Contrôle / Appels système}

    \begin{tabularx}{\columnwidth}{lX}
        \keys{\ctrl + C} & interrompre (\emph{sigint})         \\
        \keys{\ctrl + Z} & suspendre (\emph{sigstp})           \\
        \keys{\ctrl + L} & effacer l'écran                     \\
        \keys{\ctrl + D} & envoyer \emph{EOF} (ferme le shell) \\
    \end{tabularx}

    \subsection*{Casse / mots}

    \begin{tabularx}{\columnwidth}{lX}
        \keys{\Alt + U} & majuscules (mot après)        \\
        \keys{\Alt + L} & minuscules (mot après)        \\
        \keys{\Alt + C} & capitalise la première lettre \\
    \end{tabularx}

    \subsection*{Historique / recherche}

    \begin{tabularx}{\columnwidth}{lX}
        \keys{\ctrl + R} & recherche historique                \\
        \keys{\ctrl + G} & annuler la recherche                \\
        \keys{\ctrl + P} & commande précédente                 \\
        \keys{\ctrl + N} & commande suivante                   \\
        \keys{\ctrl + S} & pause (débloquer: \keys{\ctrl + Q}) \\
        \keys{\Alt + <}  & début de l'historique               \\
        \keys{\Alt + >}  & fin de l'historique                 \\
        \keys{\Alt + P}  & précédent dans l'historique         \\
        \keys{\Alt + N}  & suivant dans l'historique           \\
    \end{tabularx}

    \section*{Sucres syntaxiques}

    \begin{tabbing}
        aaaaaaaaaaaa \= a \kill
        \code|~| \> identique à \codx{\$HOME} \\
        \code|.| \> répertoire courant \\
        \code|..| \> répertoire parent \\
        \code|!!| \> dernière commande \\
        \code|!n| \> commande \#n de l'historique \\
        \codx{\$?} \> code retour dernière commande \\
        \codx{\$0} \> nom du script \\
        \code{!*} \> arguments ancienne command \\
        \code{^abc^123} \> substitution dans dernière cmd. \\
        \code{`cmd`} \> obsolète, préférer \codx{\$(cmd)} \\
        \codx{\$(cmd)} \> substitution de commande \\
        \codx{\${var}} \> rendu de variable \\
        \code{{a,b,c}} \> expansion par alternatives \\
        \code{{1..5}} \> expansions de séquence \\
        \codx{\$(cmd1; cmd2)} \> sous-coquille \\
        \codx{\$(cmd \&)} \> arrière-plan (sous-coquille) \\
        \codx{\$(( expr ))} \> arithmétique \\
    \end{tabbing}

    \section*{Redirections / pipes}
    \begin{tabbing}
        aaaaaaaaa \= a \kill
        \code|[p] > [f]| \> stdout \textrightarrow{} fichier (écrase) \\
        \code|[p] >> [f]| \> stdout \textrightarrow{} fichier (ajout) \\
        \code|[p] 2> [f]| \> stderr \textrightarrow{} fichier (écrase) \\
        \code|[p] &> [f]| \> stdout+stderr \textrightarrow{} fichier (écrase) \\
        \code|[p] &>> [f]| \> stdout+stderr \textrightarrow{} fichier (ajout) \\
        \code|[p] < [f]| \> stdin depuis fichier \\
        \code|[p] << EOF| \> heredoc (jusqu'à EOF) \\
        \code|[p] <<< "s"| \> here-string (inject sur stdin) \\
        \code{[p] | [p]} \> pipe (stdout \textrightarrow{} stdin) \\
        \code{[p] |& [p]} \> pipe (stdout+stderr \textrightarrow{} stdin) \\
    \end{tabbing}

    \section*{Exécution de commandes}
    \begin{cmdblock}
        \code|source [f]| \quad exécute script dans contexte \\
        \code|./[p]| \quad exécute programme dans rép. courant \\
    \end{cmdblock}

    \section*{Information utilisateur}
    \begin{cmdblock}
        \code|uptime| \quad temps écoulé depuis reboot \\
        \code|date| \quad date et heure actuelles \\
        \code|cal| \quad calendrier du mois en cours \\
        \code|whoami| \quad nom de l'utilisateur courant \\
        \code|hostname| \quad nom de la machine \\
        \code|arch| \quad architecture matérielle \\
        \code|lsb_release -a| \quad infos distribution Linux \\
    \end{cmdblock}

    \section*{Fichiers / répertoires}

    \begin{cmdblock}
        \code|pwd| \quad chemin répertoire courant \\
        \code|ls| \quad liste les fichiers \\
        \cmdopt{\codx{-l} détails, \codx{-a} fichiers cachés,
            \codx{-h} tailles lisibles, \codx{-R} récursif,
            \codx{-t} trier par date, \codx{-r} ordre inverse,
            \codx{-Q} guillemets, \codx{-i} inode,
            \codx{-d} répertoires seulement, \codx{-S} trier par taille,
            \codx{-1} un par ligne}
        \code|cd [d]| \quad change de répertoire \\
        \code|touch [f]| \quad crée fichier vide ou met à jour date \\
        \code|mkdir [d]| \quad crée un répertoire. \\
        \cmdopt{\codx{-v} verbeux,
        \codx{-p} parents si besoin,
        \codx{-m [perms]} permissions}
        \code|rmdir [d]| \quad supprime un répertoire vide \\
        \code|rm [f]| \quad supprime un fichier. \\
        \cmdopt{\codx{-i} interactif,
            \codx{-f} force,
            \codx{-v} verbeux,
            \codx{-r} récursif, (\codx{rm -rf} dangereux)}

        \code|cp [src] [dst]| \quad copie un fichier/répertoire. \\
        \cmdopt{\codx{-a} archive (récursif + attributs),
            \codx{-i} interactif,
            \codx{-f} force,
            \codx{-v} verbeux,
            \codx{-r} récursif}
        \code|mv [src] [dst]| \quad déplace/renomme. \\
        \cmdopt{\codx{-i} interactif,
            \codx{-f} force,
            \codx{-n} ne pas écraser,
            \codx{-v} verbeux}
        \code|ln -s [src] [dst]| \quad crée un lien symbolique \\
    \end{cmdblock}

    \section*{Consultation de fichiers}

    \begin{cmdblock}
        \code|file [f]|\quad type du fichier\par
        \code|stat [f]|\quad infos détaillées\par
        \code|wc [f]|\quad compte occurences\par
        \cmdopt{\codx{-l} lignes, \codx{-w} mots, \codx{-c} caractères}
        \code{head [f]}\quad affiche premières lignes\par
        \cmdopt{\codx{-n N} N premières lignes}
        \code|tail [f]| \quad affiche les dernières lignes \\
        \cmdopt{\codx{-n N} N dernières lignes,
            \codx{-f} temps réel}
        \code|less [f]| \quad affiche page par page \\
        \code|cat [f]| \quad affiche le contenu \\
        \code|hexdump -C [f]| \quad affiche en hexadécimal \\
        \code|diff [f1] [f2]| \quad comparaison ligne par ligne \\
        \code|cmp [f1] [f2]| \quad comparaison byte par byte \\
        \code|sha256sum [f]| \quad somme de contrôle SHA256 \\
    \end{cmdblock}

    \section*{Variables d'environnement}
    \begin{tabbing}
        aaaaaaaaaaaaa \= a \kill
        \code|set| \> liste variables et fonctions \\
        \code|env| \> liste les variables d'env. \\
        \code|export VAR=[s]| \> crée/mod. variable d'env. \\
        \code|echo \$VAR| \> affiche valeur variable \\
        \codx{\$PATH} \> chemin des exécutables \\
        \codx{\$HOME} \> répertoire personnel \\
        \codx{\$SHELL} \> interpréteur de commandes \\
        \codx{\$USER} \> nom de l'utilisateur \\
        \codx{\$PS1} \> invite de commande \\
        \codx{\$PWD} \> répertoire courant \\
        \codx{\$OLDPWD} \> répertoire précédent \\
    \end{tabbing}

    \section*{Processus / jobs}
    \begin{tabbing}
        aaaaaaaaaaaaa \= a \kill
        \code|[p]&| \> exécute en arrière-plan \\
        \code|jobs| \> liste les jobs en cours \\
        \code|wait [PID]| \> attend la fin d'un job \\
        \code|disown [PID]| \> détache un job du shell \\
        \code|fg [j]| \> ramène un job au premier plan \\
        \code|bg [j]| \> relance un job en arrière-plan \\
        \code|ps aux| \> liste les processus \\
        \code|kill [pid]| \> termine un processus \\
        \code|killall [p]| \> termine tous les processus nommés \\
        \code|top| \> moniteur de processus \\
        \code|htop| \> top amélioré (si installé) \\

    \end{tabbing}

    \section*{Archives / compression}
    \begin{cmdblock}
        \code|tar -czvf archive.tar.gz [f]| \quad crée archive .tar.gz \\
        \code|tar -xzvf archive.tar.gz| \quad extrait archive .tar.gz \\
        \code|zip archive.zip [f]| \quad crée archive .zip \\
        \code|unzip archive.zip| \quad extrait archive .zip \\
        \code|tar -cJvf archive.tar.xz [f]| \quad crée archive .tar.xz \\
        \code|tar -xJvf archive.tar.xz| \quad extrait archive .tar.xz \\
        \code|7z a archive.7z [f]| \quad crée archive .7z \\
        \code|7z x archive.7z| \quad extrait archive .7z \\
        \code|bsdtar -xf [f]| \quad extrait \\
        \code|bsdtar -cf archive.tar [f]| \quad crée archive .tar \\
    \end{cmdblock}

    \section*{Permissions / utilisateurs}

    \begin{cmdblock}
        \code|chmod| \quad modifie les permissions \\
        \cmdopt{\codx{u} user, \codx{g} group, \codx{o} others, \codx{a} all,
            \codx{+} add, \codx{-} remove, \codx{=} set exactly}
        \code|chown| \quad change le propriétaire \\
        \code|chgrp| \quad change le groupe \\
        \code|umask| \quad masque de création de fichiers \\
        \code|sudo [p]| \quad exécute en superutilisateur \\
        \code|su [utilisateur]| \quad change d'utilisateur \\
        \code|who| \quad qui est connecté \\
        \code|w| \quad qui est connecté et que font-ils \\
        \code|id| \quad informations sur l'utilisateur courant \\
        \code|groups| \quad groupes de l'utilisateur courant \\
        \code|passwd| \quad change le mot de passe \\
        \code|last| \quad historique des connexions \\
    \end{cmdblock}

    \textit{bits :} 1 x, 2 w, 4 r

    \section*{Raccourcis less / more / man}
    \begin{cmdblock}
        Espace: page suivante, \keys{b}: page précédente \\
        \keys{\return}{}: ligne suivante, \keys{y}: ligne précédente \\
        \keys{/}: rechercher, \keys{n}/\keys{N}: suivant/précédent \\
        \keys{g}/\keys{G}: début/fin, \keys{q}: quitter \\
        \keys{h}: aide, \keys{m}: marque-page \\
        \keys{'}: aller au marque-page, \keys{`}: revenir \\
    \end{cmdblock}


    \section*{Recherche}
    \begin{cmdblock}
        \code|grep "s" [f]| \quad cherche une chaîne dans un fichier \\
        \cmdopt{\codx{-P} regex Perl,
            \codx{-r} récursif,
            \codx{-i} ignore la casse,
            \codx{-v} inverse (non),
            \codx{-n} numéro de ligne,
            \codx{-c} compte les lignes,
            \codx{-l} liste les fichiers,
            \codx{-w} mot entier,
            \codx{-A n} n lignes après,
            \codx{-B n} n lignes avant,
            \codx{-C n} n lignes contexte,
            \codx{-o} seulement correspondance}
        \code|rg "s" [f]| \quad comme grep, plus rapide (si installé) \\
        \code|find [d] -name "s"| \quad cherche un fichier \\
        \cmdopt{\codx{-iname} ignore la casse,
            \codx{-type f/d} fichier/répertoire,
            \codx{-exec cmd \{\} \;} exécute cmd sur chaque fichier}
        \code|locate "s"| \quad cherche un fichier (base de données) \\
        \code|which [p]| \quad chemin d'une commande \\
    \end{cmdblock}

    \section*{Remplacement}
    \begin{cmdblock}
        \code|sed 's/search/new/g' [f]| \quad remplace avec BRE regex \\
        \cmdopt{\codx{-i} modifie le fichier,
            \codx{-n} n'affiche rien par défaut,
            \codx{-E} regex étendues,
            \codx{-e} script, \codx{-r} regex étendues}
        \code|perl -pe 's/search/new/g' [f]| \quad remplace avec regex Perl \\
        \cmdopt{\codx{-i} modifie le fichier,
            \codx{-n} n'affiche rien par défaut,
            \codx{-e} script, \codx{-l} gère les nouvelles lignes}
        \code|perl -0777 -pe 's/search/new/gs' [f]| \quad remplace sur tout le fichier \\
    \end{cmdblock}

    \section*{Utilitaires}
    \begin{cmdblock}
        \code|uniq [f]| \quad filtre les lignes consécutives identiques \\
        \cmdopt{\codx{-c} compte, \codx{-d} que les doublons,
            \codx{-u} que les uniques, \codx{-i} ignore la casse}
        \code|sort [f]| \quad trie les lignes \\
        \cmdopt{\codx{-r} inverse, \codx{-n} numérique,
            \codx{-k n} clé n, \codx{-u} unique,
            \codx{-t x} séparateur x}
        \code|cut -d x -f n [f]| \quad extrait le champ n, séparé par x \\
        \code|awk '{print \$n}' [f]| \quad extrait le champ n (espace séparateur) \\
        \code|xargs [p]| \quad construit des lignes de commande \\
        \cmdopt{\codx{-n N} N arguments par commande,
            \codx{-I \{\}} remplace \{\} par l'argument}
        \code|tr 'set1' 'set2'| \quad remplace set1 par set2 \\
        \cmdopt{\codx{-d} supprime set1, \codx{-s} squeeze}
        \code|tee [f]| \quad duplique le flux vers un fichier \\
        \cmdopt{\codx{-a} ajoute au fichier, \codx{-i} ignore les signaux}
    \end{cmdblock}

    \section*{Aide}
    \begin{tabbing}
        aaaaaaaaaaaaaa \= a \kill
        \code|help [p]| \> aide interne bash \\
        \code|[p] --help| \> aide d'une commande \\
        \code|man [p]| \> manuel d'une commande \\
        \code|apropos [terme]| \> cherche dans les manuels \\
        \code|whatis [p]| \> description commande \\
    \end{tabbing}

    \emph{Sections du manuel: 1~utilisateur, 2~syscalls, 3~lib, 4~spéciaux, 5~formats, 6~jeux, 7~divers, 8~admin.} (\code{man 3 printf})

    \section*{Gestion Réseau}
    \begin{tabbing}
        aaaaaaaaaaaaa \= a \kill
        \code|ip a| \> liste adresses IP \\
        \code|ip r| \> table de routage \\
        \code|ping [hôte]| \> teste la connectivité \\
        \code|ss -tuln| \> sockets réseau \\
        \code|dig [hôte]| \> résolution DNS avancée \\
    \end{tabbing}

    \section*{Internet}
    \begin{cmdblock}
        \code|ssh user@host| \quad connexion SSH \\
        \code|scp [src] [dst]| \quad copie sécurisée \\
        \code|wget [url]| \quad télécharge un fichier \\
        \cmdopt{\codx{-O f} fichier de sortie,
            \codx{-c} reprendre, \codx{-q} silencieux,
            \codx{--limit-rate=100k} limite le débit}
        \code|curl [url]| \quad transfère une URL \\
        \cmdopt{\codx{-o f} fichier de sortie,
            \codx{-O} nomme comme l'original,
            \codx{-L} suivre redirections,
            \codx{-I} en-têtes seulement,
            \codx{-d data} données POST,
            \codx{-H 'Hdr: v'} ajoute un en-tête}
        \code|lynx [url]| \quad navigateur web en console \\
    \end{cmdblock}

    \section*{Git}
    \begin{tabbing}
        aaaaaaaaaaaaaaaa \= a \kill
        \code|git init| \> initialise un dépôt git \\
        \code|git status| \> statut du dépôt \\
        %\code|git remote add origin [url]| \> ajoute une origine \\
        \code|git remote -v| \> liste les origines \\
        \code|git fetch| récupère les changements distants \\
        \code|git pull| \> pull + merge \\
        \code|git push| \> pousse les changements \\
        \code|git clone [url]| \> clone un dépôt distant \\
        \code|git branch| \> liste les branches \\
        \code|git branch [br]| \> crée une branche \\
        \code|git checkout [br]| \> change de branche \\
        \code|git commit -m "s"| \> crée un commit \\
        \code|git add [f]| \> ajoute un fichier au commit \\
        \code|git log| \> historique des commits \\
        \cmdopt{\codx{--oneline} résumé, \codx{--graph} graphe,
            \codx{--stat} stats, \codx{-p} diff}
    \end{tabbing}

    \section*{GCC}
    \begin{cmdblock}
        \code|gcc [f.c]| \quad compile un fichier C \\
        \cmdopt{\codx{-o f} nom de l'exécutable,
            \codx{-Wall} tous les avertissements,
            \codx{-Wextra} avertissements supplémentaires,
            \codx{-g} infos debug,
            \codx{-O0, -O1, -O2, -O3} optimisation,
            \codx{-std=c17} standard C,
            \codx{-I dir} inclure répertoire dir,
            \codx{-L dir} lier répertoire dir,
            \codx{-l lib} lier la bibliothèque lib
            \codx{-c} compiler seulement (pas linker)}
    \end{cmdblock}

    \section*{Python}
    \begin{tabbing}
        aaaaaaaaaaaaaaaaaaaaa \= a \kill
        \code|python -mvenv venv| \> crée un virtual env. \\
        \code|uv venv venv| \> même chose avec uv \\
        \code|source venv/bin/activate| \> active l'env. virtuel \\
        \code|deactivate| \> désactive l'env. virtuel \\
        \code|pip list| \> liste les paquets installés \\
        \code|pip install pkg| installe un paquet dans l'env. \\
    \end{tabbing}

    \section*{Système}

    \begin{tabbing}
        aaaaaaaaaaaaaaaaaaaaaa \= a \kill
        \code|shutdown -h now| \> arrête le système \\
        \code|shutdown -r now| \> redémarre le système \\
        \code|systemctl status [srv]| \> statut d'un service \\
        \code|systemctl start [srv]| \> démarre un service \\
        \code|systemctl stop [srv]| \> arrête un service \\
        \code|systemctl restart [srv]| \> redémarre service \\
        \code|systemctl enable [srv]| \> active service \\
        \code|systemctl disable [srv]| \> désactive service \\
        \code|journalctl -xe| \> logs système \\
        \code|timedatectl| \quad date, heure, fuseau horaire \\
        \code|hostnamectl| \quad informations sur le système \\
    \end{tabbing}

    \section*{Hardware}

    \begin{tabbing}
        aaaaaaaaaaaaaaaa \= a \kill
        \code|uname -a| \> informations système \\
        \code|lscpu| \> infos CPU \\
        \code|lsblk| \> liste les disques et partitions \\
        \code|df -h| \> espace disque \\
        \code|du -sh [f/d]| \> taille d'un fichier/répertoire \\
        \code|free -h| \> mémoire \\
    \end{tabbing}

    \section*{APT (Debian/Ubuntu)}
    \begin{tabbing}
        aaaaaaaaaaaaaaaa \= a \kill
        \code|apt update| \> met à jour liste des paquets \\
        \code|apt upgrade| \> met à jour les paquets installés \\
        \code|apt install [pkg]| \> installe un paquet \\
        \code|apt remove [pkg]| \> désinstalle un paquet \\
        \code|apt search [terme]| \> cherche un paquet \\
        \code|apt show [pkg]| \> informations sur un paquet \\
        \code|apt autoremove| \> supprime dépendances inutiles \\
        \code|dpkg -L [pkg]| \> liste fichiers d'un paquet \\
        \code|dpkg -S [f]| \> quel paquet fournit ce fichier \\
    \end{tabbing}

    \section*{Fun}

    \begin{tabbing}
        aaaaaaaaaaaaaaa \= a \kill
        \code|cowsay "s"| \> phylactere et vache \\
        \code|fortune| \> citation aléatoire  \\
        \code|sl| \> locomotive à vapeur  \\
        \code|cmatrix| \> effet Matrix  \\
        \code|toilet "s"| \> texte ASCII art  \\
        \code|lolcat| \> texte coloré  \\
    \end{tabbing}

    \section*{Substituts modernes conseillés}

    \begin{tabularx}{\columnwidth}{lX}
        lsd          & Substitut de ls             \\
        batcat       & Amélioration de cat         \\
        fd           & Substitut de find           \\
        ripgrep (rg) & Substitut de grep           \\
        sd           & Substitut de sed            \\
        hexyl        & Substitut de hexdump        \\
        broot        & Explorateur de fichiers CLI \\
        fzf          & Outil de recherche floue    \\
        git-delta    & Substitut de git diff       \\
    \end{tabularx}

\end{multicols}
\end{document}
